\section{Conventions, parameters, and the matrix target}
\label{sec:setup}

We will prove the following intermediate statement.

\begin{theorem}[Matrix construction]
\label{setup:matrix-target}
There exist a finitely generated torsion-free discrete group
\(\Gbase\), a finite integer \(N\ge1\), and a self-adjoint element
\(T\in M_N(\mathbb C[\Gbase])\), in its regular representation,
with a nonempty separated spectral subset whose spectral projection
\(p_0\) satisfies
\[
 p_0\in M_N(C_r^*(\Gbase)),\qquad
 0<\tau_{\Gbase,N}(p_0)<\frac12.
\]
\end{theorem}

Here and below a spectral subset is \emph{separated} if its distance
from the rest of the spectrum is positive. Its indicator is
continuous on the spectrum. The spectral projection therefore
belongs to the unital \(C^*\)-algebra containing the operator,
by continuous functional calculus. We also use the equivalent
contour-resolvent formula. These facts, together with polar
decomposition, are standard bounded-operator spectral theory
\cite[Chapters~VI--VII]{ReedSimon1980}.

In graph coordinates it is convenient to use right translations
\[
 (\rho_\Lambda(g)\xi)(h)=\xi(hg).
\]
The unitary \(\mathcal I\xi(h)=\xi(h^{-1})\) satisfies
\(\mathcal I\rho_\Lambda(g)\mathcal I=\lambda_\Lambda(g)\).
This identification fixes the identity vector and leaves finite
matrix indices unchanged, so it preserves all traces used here.
Thus a finite matrix sum of right translations specifies the same
kind of reduced-algebra element as in Theorem~\ref{setup:matrix-target}.

\subsection{Graphs and labels}

Edges have length one and may be traversed in either direction.
Whenever a positive orientation is specified, reverse traversal
carries the formal inverse label.
A labeled graph is \emph{folded} if, at each vertex, a prescribed
oriented label occurs on at most one outgoing edge.
A path is nonbacktracking when it never immediately reverses an
edge; a closed path is cyclically nonbacktracking if this also
holds across its chosen basepoint.
Closed walks and their repeated traversals are allowed as
presentation relators.

For a labeled group presentation we retain the formal Cayley
edges: a positive generator and its initial vertex specify an
edge even if different formal generators have the same group
value. All degree and forest assertions are about this graph.
This convention also specifies an unambiguous left action on
oriented edges. Intrinsic distances and balls in an embedded
copy are measured using its own graph edges.

All vertex Hilbert spaces use counting measure.
A biregular incidence from a \(B\)-side of degree \(D'\) to an
\(A\)-side of degree \(d\) is normalized by
\((dD')^{-1/2}\). It maps the normalized constant vector on
one side to the normalized constant vector on the other.
When using probabilities to prove a bound for that operator,
we pass explicitly to uniform probability on each side.

\subsection{The growing parameter}

Let \(\ell\) be a positive integer, eventually sufficiently large,
and set
\begin{equation}
\begin{aligned}
 d&=2^{4\ell},& D&=d^2,& s&=10\ell,& r&=7\ell,\\
 n_0&=1000,&
 H_i&=(2n_0+1)s+1+|i| \quad(i\in\mathbb Z).
\end{aligned}
\label{setup:parameters}
\end{equation}
Write \(H_0=(2n_0+1)s+1\), and put
\begin{equation}
 L_i=\frac{100H_i}{\log_2d},\qquad
 g_i=\frac{H_i}{100\log_2d}.
 \label{setup:lengths}
\end{equation}
Inequalities comparing integer walk lengths with these real
numbers have their usual literal meaning; a lower required
length can equivalently be rounded upward.
We may use the fixed bound
\begin{equation}
 H_0\le h_*\log_2d,\qquad h_*=5003,
 \label{setup:hstar}
\end{equation}
valid for every \(\ell\ge1\).

Edges in the base tower change the layer by at most one.
Consequently, for any fixed \(C<\infty\),
\begin{equation}
 |j-i|\le \frac{CH_i}{\log_2d}
 \quad\Longrightarrow\quad
 \left|\frac{H_j}{H_i}-1\right|
 \le\frac{C}{\log_2d}.
 \label{setup:nearby}
\end{equation}
The convergence to one is uniform in \(i\). Constants such as
\(C\) may be very large, but are fixed before \(\ell\) grows.

Unless stated otherwise, a term \(o(1)\) tends to zero as
\(\ell\to\infty\), uniformly over all allowed layers, blocks,
and choices already declared arbitrary.
Fixed constants may depend on fixed accuracies and earlier
parameters. They may not depend on \(\ell\) or a layer index.
Later sections record the order of the remaining choices:
geometric constants first, then nilpotence and mask scales,
then operator accuracies and finite test constants, and
finally a sufficiently large finite \(\ell\).
We do not choose different values of \(\ell\) for different layers.

\subsection{The order of choices}

Table~\ref{tab:projection-parameters} records which choices are fixed before
which others. It is a guide to the estimates proved below: each indicated
section supplies the corresponding inequality. All choices use one final
value of \(\ell\), uniformly over the infinitely many layers.

\begin{table}[htbp]
\centering
\small
\begin{tabular}{@{}>{\raggedright\arraybackslash}p{0.25\textwidth}>{\raggedright\arraybackslash}p{0.38\textwidth}>{\raggedright\arraybackslash}p{0.29\textwidth}@{}}
\toprule
Choice & Required conclusion & Proof location \\
\midrule
Geometric constants \(C_*,C_\bullet\) &
One ball contains each earlier overlap; the unused edges form a forest &
Section~\ref{geo:section} \\[3pt]
Nilpotence and scale constants \(m,E_0,A_*,\kappa\) &
Short words are detected; both endpoint regimes have small weighted trace &
Section~\ref{nil:section} \\[3pt]
Accuracy \(\eta\), coefficient cutoffs and finite lists &
Uniform approximation in the full regular representation with a finite weight alphabet &
Section~\ref{sam:section}; Equation~\eqref{spec:eta-choice} \\[3pt]
Finite-test and moment constants \(C_1,C_2\) &
The sum of all voltage and operator failure bounds is less than one &
Equation~\eqref{sam:C2-order} and Theorem~\ref{sam:realization} \\[3pt]
One sufficiently large \(\ell\) &
Deletion loss, weighted trace, and forest error meet all spectral tolerances &
Theorem~\ref{spec:matrix-projection} \\
\bottomrule
\end{tabular}
\caption{Parameter ledger for the projection construction. Constants in a row
may depend on earlier fixed choices; only the final row increases \(\ell\).
The matrix dimension \(N\) is a prescribed function of \(\ell\), and its
unnormalized contribution is retained in the trace estimates.}
\label{tab:projection-parameters}
\end{table}
